\documentclass[11pt, a4paper]{article}
\usepackage[margin=2.5cm]{geometry}
\usepackage{parskip}
\usepackage{booktabs}
\usepackage{hyperref}
\usepackage{amsmath}
\usepackage{xcolor}
\definecolor{noteblue}{RGB}{30, 90, 160}

\title{\textbf{Story Notes: Climate and Influenza A HA Evolution}\\
\large Key findings, methods sketch, and open questions}
\author{Nina Pust}
\date{\today}

\begin{document}
\maketitle

% ─────────────────────────────────────────────────────────────────────────────
\section*{Pipeline overview}

The analysis takes GISAID Influenza A sequences and matched ERA5-Land climate
data and asks whether environmental conditions predict molecular evolution of
the haemagglutinin (HA) surface protein.

\textbf{Sequences and metadata.}
Sequences for H1N1 (H1, N1) and H3N2 (H3, N2) surface proteins plus PA
(internal protein, negative control) were exported from GISAID EpiFlu.
Metadata were geocoded via Nominatim and monthly ERA5-Land temperature and
relative humidity (derived from dewpoint via Magnus formula) were attached to
each sequence by collection location and date.

\textbf{Phylogenetics.}
Sequences were aligned with MAFFT, codon-aligned with MACSE (pairwise
batch approach), and initial trees built with MAPLE.
Two-round pruning removed duplicate/identical sequences (zero-length branches)
and outliers (unusually long branches).
Trees were dated with TreeTime to obtain a timetree with branch lengths in
years.

\textbf{Evolutionary rates.}
Two parallel approaches were used:
\begin{enumerate}
  \item \textit{MAPLE/timetree branch rate:} total substitutions per year on
        each terminal branch, as estimated directly from the timetree
        (\texttt{log\_rate}).
  \item \textit{Ancestral parsimony reconstruction:} for each codon in the
        alignment, ancestral states were reconstructed at internal nodes via
        maximum parsimony (\texttt{phangorn::ancestral.pars}, MPR).
        Non-synonymous (NS) and synonymous (S) changes on each terminal branch
        were counted and divided by timetree branch length to give per-year
        rates. This was applied to three regions: the 7-codon antigenic ridge,
        the full epitope ($\sim$103 codons), and the whole HA gene
        ($\sim$579 codons). The approach avoids the dS denominator problem that
        makes yn00-based $\omega$ unreliable on short or sparse alignments.
  \item \textit{yn00 $\omega$ (dN/dS):} pairwise yn00 estimates mapped onto
        the tree via NNLS and FastME. Used for epitope-level analyses but
        largely superseded by ancestral parsimony; results were null after
        applying appropriate rate floors.
\end{enumerate}

\textbf{Climate predictors.}
All predictors were scaled (mean 0, SD 1):
temperature (\texttt{scaled\_temp}),
relative humidity (\texttt{scaled\_humidity}),
$\Delta$temperature and $\Delta$humidity (change from previous sampling
interval), sampling density (to absorb seasonal confounding), and population
density (not yet fully integrated; largely NA).

\textbf{Statistical models.}
Phylogenetic regression via \texttt{phylolm} (Brownian motion covariance) was
used to test climate predictors against log-transformed rates, controlling for
year as a factor and for phylogenetic non-independence.
Sampling density was included as a covariate to absorb the seasonal
flu/sampling winter peak, so surviving temperature signals reflect
environmental effects beyond seasonality.
BH-FDR correction was applied across all predictor $\times$ response
combinations within each gene.
A Wilcoxon rank-sum test with rank-biserial effect size was used to test
whether specific recurrent mutations (${>}5$ independent occurrences) arise
preferentially under particular climate conditions.

% ─────────────────────────────────────────────────────────────────────────────
\section*{Key findings}

\subsection*{1. Faster HA evolution in colder climates}

This is the central finding and it is \textbf{consistent across methods,
genes, and genomic regions}.

\begin{itemize}
  \item \textbf{MAPLE/timetree branch rate (total substitutions/year):}
        Temperature has a significant negative effect on \texttt{log\_rate}
        in both H1N1 (estimate $= -0.047$, BH-FDR $p = 3 \times 10^{-4}$)
        and H3N2 (estimate $= -0.030$, BH-FDR $p = 0.016$) after controlling
        for sampling density, year, and humidity.

  \item \textbf{Ancestral parsimony — full gene NS rate:}
        H1N1 estimate $= -0.059$ (FDR $p = 4 \times 10^{-6}$);
        H3N2 estimate $= -0.041$ (FDR $p = 4.9 \times 10^{-3}$).

  \item \textbf{Ancestral parsimony — full gene S rate:}
        H1N1 estimate $= -0.099$ (FDR $p = 4 \times 10^{-6}$);
        H3N2 estimate $= -0.035$ (FDR $p = 6.4 \times 10^{-3}$).
        The effect on synonymous rates — which are largely neutral — is
        noteworthy: it points to a \textit{mutational} rather than purely
        selective mechanism.

  \item \textbf{Epitope NS rate:}
        H1N1 estimate $= -0.106$ (FDR $p = 6.5 \times 10^{-5}$);
        H3N2 (marginal, FDR $p \approx 0.10$).

  \item \textbf{Ridge NS rate:}
        H1N1 estimate $= -0.186$ (FDR $p = 0.014$). Weaker signal,
        consistent direction.
\end{itemize}

\textbf{Interpretation.}
The negative temperature coefficient means lower temperatures are associated
with higher per-branch evolutionary rates.
This aligns with the winter seasonality of influenza: cold conditions
coincide with peak transmission, faster replication, and potentially elevated
polymerase error rates.
The fact that both NS and S rates are affected, and that the effect survives
after controlling for sampling density (which absorbs the seasonal flu peak),
supports a genuine environmental effect on the HA mutation rate rather than a
sampling artefact.

\textbf{Sampling density as a control.}
Sampling density is positively associated with rates in several analyses
(particularly H3N2, FDR $p < 10^{-4}$), consistent with denser sampling
capturing more short-branch diversity. Its inclusion is necessary; temperature
results should not be interpreted from models that omit it.

% ─────────────────────────────────────────────────────────────────────────────
\subsection*{2. Jackknife subsampling confirms signal is distributed across
the dataset}

A block jackknife was run for H1 and H3: 1{,}000 random subsamples of 1{,}000
sequences each, with \texttt{phylolm} (log\_rate $\sim$ scaled\_temp +
year) fitted independently to each subsample.

\textbf{Results.}
\begin{itemize}
  \item \textbf{H1:} 150/1{,}000 subsamples were significant at $p < 0.001$
        ($***$), with temperature coefficient estimates distributed
        consistently around a negative value, matching the full-dataset
        result (estimate $= -0.047$, BH-FDR $p = 3 \times 10^{-4}$).
  \item \textbf{H3:} Only 1/1{,}000 subsamples reached $p < 0.001$, consistent
        with the weaker full-dataset signal (BH-FDR $p = 0.016$).
        H3 behaves exactly as expected for a real but marginal effect that
        requires the full sample size to be reliably detectable.
\end{itemize}

\textbf{Are the significant subsamples driven by specific sequences?}
Enrichment analysis (ratio of sequence frequency in $***$ vs.\ ns subsamples)
identified 9 sequences with ratio $> 2$ for H1. These were not biological
drivers:
\begin{itemize}
  \item Four of the nine (EPI\_ISL\_70358/59/62/63) are from the same outbreak
        cluster (Bosnia, January 2009, identical temperature $-2.46^\circ$C and
        humidity 84.1\%). They co-occur in subsamples by construction and
        collectively pull the coefficient when sampled together, but represent
        at most one or two independent observations.
  \item The remaining five span diverse geographies (Georgia, Singapore,
        Costa Rica, Brazil, Russia) with no shared climate profile.
  \item For H3, the single $***$ subsample causes the enrichment metric to
        collapse entirely --- every sequence in that one subsample appears
        10$\times$ enriched, rendering the metric meaningless.
\end{itemize}

\textbf{Conclusion.}
If the signal were artefactual or driven by outlier sequences, the same
sequences would recur consistently in significant subsamples, producing a
stable enrichment pattern. Instead, significant subsamples are largely
interchangeable in their sequence composition.
The jackknife therefore supports the interpretation that \textbf{the
temperature effect on HA evolutionary rate is real and distributed across the
full dataset}, rather than being propped up by a small number of influential
observations.

% ─────────────────────────────────────────────────────────────────────────────
\subsection*{3. No specific recurrent mutation is climate-associated}

Wilcoxon rank-sum tests were applied to all recurrent NS mutations
($n > 5$ independent terminal-branch occurrences) in the antigenic ridge and
full epitope, testing association with temperature, humidity,
$\Delta$temperature, and $\Delta$humidity.

After BH-FDR correction, \textbf{no mutation in the ridge survived} in either
H1N1 or H3N2. In the epitope, a single H3N2 mutation at residue 47
(H47R, $n = 9$, rank-biserial $r = -0.73$, FDR $p = 0.044$) was significant
for temperature; this is a single finding among hundreds of tests and should
be treated cautiously pending replication.

\textbf{Interpretation.}
The climate effect on HA evolution appears to operate on the \textit{overall
rate} rather than driving specific adaptive substitutions at identifiable
positions. This is consistent with a mutational (polymerase fidelity) rather
than a selective (immune-driven fixation) mechanism.

% ─────────────────────────────────────────────────────────────────────────────
\subsection*{4. Bayesian phylogenetic regression not reliable on subsamples}

Both INLA and \texttt{brms} were tested as Bayesian alternatives to
\texttt{phylolm}.

\textbf{INLA} was found to be sensitive to the random seed used for subset
sampling, with different seeds producing qualitatively different conclusions.
This likely reflects the large dataset size and the complexity of the spatial
and phylogenetic random effects. INLA results are excluded entirely.

\textbf{brms} (7 seeds, 1{,}000-sequence subsamples) showed the same
instability:
\begin{itemize}
  \item \textbf{H1:} Only seed 42 produced a credible interval excluding zero
        (mean $= -0.067$, 95\% CI $[-0.115, -0.018]$). The remaining six seeds
        gave estimates ranging from $-0.046$ to $+0.004$, none significant.
  \item \textbf{H3:} No seed reached significance; estimates ranged from
        $-0.048$ to $+0.022$ and changed sign across seeds.
\end{itemize}

The instability reflects insufficient power at $n = 1{,}000$ for a Bayesian
phylogenetic model to reliably recover a signal that \texttt{phylolm} detects
only at $n \approx 9{,}000$--$13{,}000$ sequences (H1 BH-FDR $p = 3 \times
10^{-4}$, H3 BH-FDR $p = 0.016$). A full-dataset \texttt{brms} run is in
progress; results pending.

\texttt{phylolm} remains the primary analysis throughout.

% ─────────────────────────────────────────────────────────────────────────────
\subsection*{5. Additional notes}

\begin{itemize}
  \item \textbf{Ridge $\omega$ artefact.}
        Early analyses showed an apparent humidity $\rightarrow$ ridge $\omega$
        association. This was traced to a dS denominator problem: with only 7
        codons, $\approx 96\%$ of ridge dS values are effectively zero, making
        $\omega = \text{dN}/\text{dS}$ undefined or inflated.
        The result disappears when ancestral parsimony rates are used instead.
        A reminder that $\omega$ is not interpretable on very short alignments.

  \item \textbf{Null result for yn00 $\omega$ after rate floors.}
        After applying log-rate floors to remove implausibly small dN and dS
        estimates (bimodal gap in distributions), all epitope-level $\omega$
        phylolm associations became non-significant (except the ridge artefact
        above). The ancestral parsimony approach is therefore preferred for all
        substitution-rate analyses.

  \item \textbf{PA negative control.}
        PA (polymerase acidic, internal protein) serves as a negative control.
        PA should show no climate signal if the HA findings reflect
        surface-protein-specific biology. Results for PA were not assessed in
        the epitope/ridge pipeline (surface-protein-specific) but should be
        checked in the main phylolm results before finalising.

  \item \textbf{Geographic confounding.}
        Sequences are dominated by temperate Northern Hemisphere sampling
        (US, Europe). Temperature is partly a proxy for geography/latitude.
        Phylolm corrects for phylogenetic clustering but geographic and
        phylogenetic structure are correlated in flu. This should be
        acknowledged as a limitation.

  \item \textbf{Population density.}
        The \texttt{scaled\_pop\_density} predictor is largely NA in the
        current dataset (source not yet integrated). A significant negative
        effect on H1N1 ridge NS rate was observed in the data that do have
        values; interpret cautiously.
\end{itemize}

% ─────────────────────────────────────────────────────────────────────────────
\section*{Open questions / next steps}

\begin{itemize}
  \item Check PA negative control in main phylolm results.
  \item Verify H47R (epitope, H3N2) in independent dataset or with larger $n$.
  \item Consider latitude or hemisphere as an explicit covariate to
        disentangle temperature from geography.
  \item Integrate population density data to complete the predictor set.
  \item Investigate whether the S-rate temperature effect differs between
        coding regions (would help distinguish polymerase fidelity vs.
        replication rate hypotheses).
\end{itemize}

\end{document}
